%% ============================================================
%%  LUCRARE DE LICENȚĂ
%%  Universitatea Politehnica Timișoara, 2026
%%  Compilare: pdflatex licenta_gorbenco.tex  (x2)
%% ============================================================

\documentclass[12pt, twoside, a4paper]{report}

%% ---- Encoding & font (Arial echivalent) ----
\usepackage[utf8]{inputenc}
\usepackage[T1]{fontenc}
\usepackage{helvet}
\renewcommand\familydefault{\sfdefault}

%% ---- Limbă ----
\usepackage[romanian]{babel}

%% ---- Margini (Mirror Margins, conform ghid UPT) ----
\usepackage[
  a4paper, twoside,
  inner=2cm, outer=2cm,
  top=2.5cm, bottom=2cm,
  includehead, headsep=5mm,
  headheight=38pt
]{geometry}

%% ---- Spațiere rânduri 1,15 ----
\usepackage{setspace}
\setstretch{1.15}

%% ---- Indent paragrafe 1,5 cm, inclusiv primul ----
\setlength{\parindent}{1.5cm}
\usepackage{indentfirst}
\setlength{\parskip}{0pt}

%% ---- Antet și subsol ----
\usepackage{fancyhdr}
\fancyhf{}
\fancyhead[L]{%
  \fontsize{8}{10}\selectfont
  Universitatea Politehnica Timișoara\\
  Ingineria Sistemelor, 2025--2026\\
  Vidac Denisa \quad\quad Sistem Smart Home Local-First%
}
\fancyfoot[C]{\thepage}
\renewcommand{\headrulewidth}{0.4pt}

%% ---- Tabele ----
\usepackage{booktabs}
\usepackage{array}
\usepackage{longtable}
\usepackage{multirow}
\usepackage{caption}
\captionsetup[table]{
  labelfont=bf, labelsep=period,
  justification=centering,
  position=above, skip=4pt
}
\captionsetup[figure]{
  labelfont=bf, labelsep=period,
  justification=centering,
  position=below, skip=4pt
}

%% ---- Grafice ----
\usepackage{graphicx}
\usepackage{float}

%% ---- Verbatim (diagrame ASCII) ----
\usepackage{fancyvrb}

%% ---- Hyperlink-uri (fără chenare colorate) ----
\usepackage[hidelinks]{hyperref}
\hypersetup{
  pdftitle  = {Sistem Smart Home Local-First cu Detecție de Anomalii
               Comportamentale bazată pe Procesare Edge},
  pdfauthor = {Vidac Denisa},
  pdfsubject= {Lucrare de licență, UPT 2026}
}

%% ---- Micro-tipografie ----
\usepackage{microtype}


%% ---- Formatare titluri capitole / secțiuni ----
\usepackage{titlesec}

%  Capitol: "Capitolul X" pe rând separat, titlu dedesubt
\titleformat{\chapter}[display]
  {\normalfont\large\bfseries}
  {\chaptertitlename\ \thechapter}
  {8pt}
  {\large\bfseries}
\titlespacing*{\chapter}{0pt}{0pt}{20pt}

%  Secțiune
\titleformat{\section}
  {\normalfont\normalsize\bfseries}
  {\thesection}{0.6em}{}
\titlespacing*{\section}{0pt}{14pt}{4pt}

%  Subsecțiune
\titleformat{\subsection}
  {\normalfont\normalsize\bfseries}
  {\thesubsection}{0.6em}{}
\titlespacing*{\subsection}{0pt}{10pt}{3pt}

%  Sub-subsecțiune
\titleformat{\subsubsection}
  {\normalfont\normalsize\bfseries}
  {\thesubsubsection}{0.6em}{}
\titlespacing*{\subsubsection}{0pt}{8pt}{2pt}

%% ---- Adâncime numerotare TOC ----
\setcounter{tocdepth}{3}
\setcounter{secnumdepth}{3}

%% ================================================================
\begin{document}
%% ================================================================


%% ============================================================
%%  COPERTA
%% ============================================================
\begin{titlepage}
\thispagestyle{empty}
\pagenumbering{arabic}   % numărătoarea începe de la copertă (p.1)
                          % dar numărul e afișat doar de la Introducere

\begin{center}

\vspace*{0.3cm}
{\fontsize{11}{13}\selectfont Universitatea Politehnica Timișoara}\\[2pt]
{\fontsize{11}{13}\selectfont Facultatea de Automatică și Calculatoare}\\[2pt]
{\fontsize{11}{13}\selectfont Ingineria Sistemelor}\\[2pt]
{\fontsize{11}{13}\selectfont Anul universitar 2025--2026}

\vfill

{\fontsize{20}{24}\bfseries\MakeUppercase{Sistem Smart Home Local-First
cu Detecție de Anomalii \\ Comportamentale bazată pe Procesare Edge}\par}

\vfill

\end{center}

\begin{flushleft}
{\fontsize{14}{17}\bfseries Candidat: Vidac Denisa\par}
\vspace{6pt}
{\fontsize{14}{17}\bfseries Coordonator științific: Conf./Prof. dr. ing.
\underline{\hspace{6cm}}\par}
\end{flushleft}

\vspace{2cm}

\begin{center}
{\fontsize{14}{17}\selectfont Sesiunea: Iunie 2026}
\end{center}

\end{titlepage}


%% ============================================================
%%  REZUMAT
%% ============================================================
\newpage
\thispagestyle{empty}
\vspace*{1cm}

\begin{center}
{\fontsize{20}{24}\bfseries\MakeUppercase{Rezumat}\par}
\end{center}

\vspace{1cm}

Lucrarea de față prezintă proiectarea și implementarea unui sistem de automatizare
rezidențială (\textit{smart home}) local-first, care funcționează în mod complet independent
față de orice infrastructură cloud externă. Sistemul este construit pe o arhitectură în
cinci straturi, utilizând microcontrolere ESP32 WROOM-32 pentru colectarea datelor de
la senzori (temperatură, umiditate, gaz, mișcare, lumină) și controlul actuatoarelor
(releu de putere, servomotor, buzzer), un modul ESP32-CAM pentru supravegherea video,
un gateway Raspberry Pi 3B+ care găzduiește broker-ul MQTT Mosquitto, backend-ul
NestJS, baza de date TimescaleDB și motorul de detecție a anomaliilor, și o aplicație
mobilă React Native/Expo pentru interfața utilizatorului.

Contribuția principală a lucrării constă în integrarea unui mecanism de detecție a
anomaliilor comportamentale bazat pe algoritmul Isolation Forest, rulat exclusiv pe
Raspberry Pi, care identifică devierile de la tiparele normale fără a transmite datele
personale în afara rețelei locale. Sistemul demonstrează că o soluție IoT de nivel
competitiv poate fi construită integral cu hardware accesibil (sub 150 EUR) și
tehnologii open-source, oferind o alternativă viabilă la ecosistemele comerciale
Google Home, Amazon Alexa sau Tuya.

\textbf{Cuvinte cheie:} IoT, smart home, local-first, edge computing, ESP32,
Raspberry Pi, MQTT, NestJS, React Native, Isolation Forest, detecție anomalii.


%% ============================================================
%%  CUPRINS
%% ============================================================
\newpage
\thispagestyle{empty}
\tableofcontents


%% ================================================================
%%  CONȚINUT PRINCIPAL – antetul apare de aici
%% ================================================================
\cleardoublepage
\pagestyle{fancy}


%% ============================================================
%%  CAPITOLUL 1 — INTRODUCERE
%% ============================================================
\chapter{Introducere}

\section{Contextul și Motivația Proiectului}

Internetul Lucrurilor (IoT --- \textit{Internet of Things}) reprezintă una dintre cele
mai semnificative transformări tehnologice ale ultimului deceniu. Conform rapoartelor
Statista și IoT Analytics, numărul dispozitivelor IoT conectate la nivel global a depășit
pragul de 15 miliarde în anul 2024, cu o prognoză de creștere spre 30 de miliarde până în
2030. În cadrul acestui ecosistem în expansiune rapidă, segmentul casei inteligente
(\textit{smart home}) ocupă o poziție de prim rang, atât din perspectiva adoptării
consumatorilor, cât și din cea a volumului de investiții al marilor corporații tehnologice.

Piața globală a sistemelor smart home a fost evaluată la aproximativ 110 miliarde de
dolari în 2024 și se estimează că va atinge 240 de miliarde de dolari până în 2030, cu o
rată de creștere anuală compusă (CAGR) de aproximativ 13,5\%. Această creștere este
impulsionată de accesibilizarea dispozitivelor senzoriale, reducerea costurilor
microcontrolerelor și creșterea penetrării rețelelor Wi-Fi în gospodăriile moderne.
Soluțiile comerciale dominante --- precum Google Home, Amazon Alexa și ecosistemul
Tuya --- au democratizat accesul la automatizarea casei, oferind interfețe intuitive și o
gamă largă de dispozitive compatibile.

Cu toate acestea, analiza critică a acestor soluții comerciale relevă o serie de
limitări structurale care ridică întrebări fundamentale legate de confidențialitate,
fiabilitate și autonomie tehnologică. Marea majoritate a ecosistemelor smart home
comerciale sunt construite în jurul unui model arhitectural \textit{cloud-centric}, în
care datele colectate de senzorii din locuință sunt transmise și stocate pe servere
externe, gestionate de companii terțe. Astfel, un senzor de temperatură amplasat în
dormitorul utilizatorului transmite permanent date către infrastructura cloud a
producătorului, iar un dispozitiv PIR (\textit{Passive Infrared}) pentru detectarea
mișcării furnizează implicit informații detaliate despre comportamentul și prezența
persoanelor în locuință.

Problemele asociate acestui model sunt multiple și de natură diversă. În primul rând,
există o dependență critică de disponibilitatea internetului și a serviciilor cloud ale
furnizorului. Numeroase cazuri documentate demonstrează că sistemele smart home bazate
pe cloud devin complet nefuncționale în momentul în care furnizorul decide să întrerupă
suportul pentru produse mai vechi sau când serverele de backend sunt temporar indisponibile.
Cazul emblematic al platformei Wink, care a trecut brusc la un model prin abonament în
2020 lăsând utilizatorii fără acces la propriile dispozitive, sau decizia Google de a
discontinua suportul pentru Nest Secure, ilustrează fragilitatea acestor dependențe.
În al doilea rând, modelul cloud-centric implică transmiterea continuă a datelor personale
sensibile --- rutine comportamentale, orare de activitate, prezența sau absența persoanelor
în locuință --- către servere externe, creând profile detaliate ale utilizatorilor cu
implicații profunde pentru confidențialitate. Regulamentul General privind Protecția
Datelor (GDPR) din Uniunea Europeană a adresat parțial aceste preocupări, însă mecanismele
de aplicare rămân limitate în fața unor companii cu infrastructuri globale complexe.

În acest context, se identifică o oportunitate clară pentru proiectarea și implementarea
unui sistem smart home alternativ, fundamentat pe principii arhitecturale diferite:
localitate completă a datelor, procesare inteligentă la marginea rețelei
(\textit{edge computing}) și independență față de orice infrastructură cloud. Progresele
recente în domeniul microcontrolerelor de putere redusă --- în special familia ESP32 de
la Espressif Systems --- și disponibilitatea unor platforme de calcul la nivel de placă
unică, precum Raspberry Pi, fac posibilă implementarea unui sistem complet funcțional
la un cost de hardware semnificativ mai mic decât cel al soluțiilor comerciale echivalente.

\section{Obiectivele Lucrării}

Pornind de la analiza contextului și a problemelor identificate în secțiunea
anterioară, prezenta lucrare urmărește realizarea unui set coerent de obiective,
structurate pe mai multe planuri ale proiectării și implementării sistemului.

\textbf{Primul obiectiv major} constă în proiectarea și implementarea unui sistem
smart home complet funcțional, bazat exclusiv pe infrastructură locală, fără nicio
dependență de servicii externe sau conectivitate la internet pentru funcționarea de bază.
Sistemul trebuie să asigure monitorizarea în timp real a parametrilor de mediu
(temperatură, umiditate, concentrație de gaze, nivelul luminii, detectarea mișcării),
controlul actuatoarelor (releu de putere, motor servo, buzzer de alertă) și gestionarea
unui nod de supraveghere video, toate acestea operând exclusiv în cadrul rețelei locale
a utilizatorului.

\textbf{Al doilea obiectiv} vizează integrarea unui mecanism de detecție a anomaliilor
comportamentale bazat pe tehnici de machine learning, rulat în întregime pe dispozitivul
gateway (Raspberry Pi 3B+). Spre deosebire de sistemele clasice de alertă bazate pe
praguri fixe (\textit{threshold-based}), abordarea propusă utilizează algoritmul
Isolation Forest \cite{liu2008} pentru a învăța tiparele normale de comportament ale
utilizatorului și pentru a detecta deviațiile semnificative de la aceste tipare. Această
capacitate conferă sistemului un grad de inteligență adaptivă, reducând rata de alerte
fals-pozitive caracteristice sistemelor bazate pe reguli rigide.

\textbf{Al treilea obiectiv} se referă la realizarea unei aplicații mobile moderne,
disponibile pe platformele iOS și Android, care să ofere utilizatorului o interfață
intuitivă pentru monitorizarea stării sistemului, vizualizarea istoricului datelor,
gestionarea alertelor și controlul de la distanță al actuatoarelor. Aplicația trebuie
să beneficieze de actualizări în timp real prin intermediul protocolului WebSocket,
eliminând necesitatea interogării periodice (\textit{polling}) a serverului și asigurând
o experiență de utilizare fluidă și reactivă.

\textbf{Al patrulea obiectiv} constă în demonstrarea că un sistem IoT de nivel
competitiv poate fi construit integral fără dependență de cloud, utilizând exclusiv
tehnologii open-source și hardware comercial accesibil. Această demonstrație are valoare
atât din perspectiva practică, cât și din cea academică, contribuind la literatura de
specialitate în domeniul arhitecturilor IoT locale și al procesării edge.

\textbf{Al cincilea obiectiv}, cu caracter pragmatic, stabilește un buget maxim de
hardware de 150 EUR pentru întreaga infrastructură fizică a sistemului, incluzând
microcontrolerele ESP32, senzorii, actuatorii, cablajul și placa Raspberry Pi.
Respectarea acestei constrângeri de cost demonstrează accesibilitatea soluției pentru
publicul larg și competitivitatea sa față de ecosistemele comerciale cu costuri lunare
recurente.

\section{Contribuția Proprie}

Lucrarea de față reprezintă o contribuție originală pe mai multe planuri tehnice
interconectate, autorul proiectând și implementând toate componentele sistemului de la
zero, fără utilizarea de soluții preconstruite sau platforme \textit{no-code}.

\textbf{Firmware-ul pentru nodurile ESP32} a fost dezvoltat în MicroPython, oferind
o soluție echilibrată între performanță și productivitate pentru gestionarea senzorilor
și actuatoarelor. Firmware-ul implementează conectivitatea MQTT cu autentificare și
reconectare automată, rutine de citire pentru senzorii DHT22 (temperatură/umiditate),
MQ-2 (gaze combustibile), PIR (mișcare) și BH1750 (luminanță), precum și controlul
actuatoarelor (releu, servo motor, buzzer).

\textbf{Backend-ul NestJS} constituie stratul de orchestrare al întregului sistem,
expunând un API REST securizat cu autentificare JWT și un gateway WebSocket pentru
transmiterea evenimentelor în timp real către aplicația mobilă. Arhitectura modulară a
aplicației NestJS --- cu module separate pentru autentificare, gestionarea senzorilor,
controlul actuatoarelor și alertare --- permite extensibilitatea și mentenabilitatea
codului. Integrarea cu TimescaleDB prin TypeORM și cu broker-ul MQTT prin biblioteca
\texttt{mqtt.js} realizează bridging-ul între protocoalele de comunicație IoT și
interfețele REST/WebSocket moderne.

\textbf{Pipeline-ul de detecție a anomaliilor} reprezintă componenta de machine
learning a sistemului, implementată în Python pe Raspberry Pi 3B+. Algoritmul
Isolation Forest, implementat prin biblioteca \texttt{scikit-learn}, este antrenat pe
date istorice de comportament ale utilizatorului și rulat periodic pentru a evalua
normalitatea citirilor curente ale senzorilor. Contribuția proprie include proiectarea
schemei de antrenare incrementală, normalizarea datelor multidimensionale și integrarea
rezultatelor algoritmului cu sistemul de alertare prin MQTT.

\textbf{Aplicația mobilă React Native/Expo} oferă o interfață utilizator modernă,
cu o temă vizuală caldă bazată pe paleta de culori stone (gri-bej) și amber (chihlimbar),
adaptată pentru utilizare nocturnă. Aplicația implementează autentificare JWT cu
persistență locală, vizualizare în timp real a datelor senzorilor prin WebSocket,
grafice de evoluție temporală și un sistem de notificări pentru alertele detectate de
algoritmul ML.

\textbf{Infrastructura locală completă} --- incluzând broker-ul Mosquitto MQTT
configurat cu TLS, baza de date TimescaleDB cu hypertable-uri optimizate pentru serii
temporale și configurarea rețelei locale --- constituie în ansamblu o soluție integrată
care demonstrează fezabilitatea paradigmei local-first pentru aplicații IoT de producție.

\section{Structura Lucrării}

Lucrarea este organizată în șase capitole principale, fiecare abordând un aspect
distinct al proiectului, de la fundamentele teoretice la detaliile de implementare și
rezultatele experimentale.

\textbf{Capitolul 2 --- Arhitectura Sistemului} prezintă viziunea arhitecturală de
ansamblu a proiectului, detaliind principiile local-first și privacy-by-design care au
ghidat toate deciziile de proiectare. Sunt descrise în detaliu cele cinci straturi ale
arhitecturii, împreună cu fluxurile de date și justificările pentru alegerea fiecărei
tehnologii componente. De asemenea, sunt analizate considerentele de securitate ale
sistemului.

\textbf{Capitolul 3 --- Implementarea Hardware} acoperă în detaliu componentele
fizice ale sistemului: microcontrolerele ESP32 WROOM-32 și ESP32-CAM, senzorii
utilizați (DHT22, MQ-2, PIR, BH1750) și actuatorii (releu, servo motor, buzzer).
Sunt prezentate schemele de conexiuni, considerentele de alimentare electrică și
configurarea rețelei fizice de noduri senzoriale.

\textbf{Capitolul 4 --- Implementarea Software} descrie în detaliu toate componentele
software ale sistemului: firmware-ul MicroPython pentru ESP32, configurarea broker-ului
Mosquitto și a bazei de date TimescaleDB pe Raspberry Pi, arhitectura backend-ului NestJS
și implementarea aplicației mobile React Native/Expo.

\textbf{Capitolul 5 --- Detecția de Anomalii prin Machine Learning} este dedicat
componentei de inteligență artificială a sistemului. Sunt prezentate fundamentele
teoretice ale algoritmului Isolation Forest, procesul de colectare și preprocesare a
datelor de antrenare, rezultatele antrenării și evaluarea performanțelor algoritmului
pe date reale colectate de sistemul implementat.

\textbf{Capitolul 6 --- Testare și Evaluare} prezintă metodologia de testare aplicată
sistemului, rezultatele testelor funcționale și de performanță, analiza consumului de
resurse pe Raspberry Pi și o evaluare comparativă față de soluțiile comerciale existente.
Sunt discutate limitările sistemului și direcțiile de dezvoltare ulterioară.


%% ============================================================
%%  CAPITOLUL 2 — ARHITECTURA SISTEMULUI
%% ============================================================
\chapter{Arhitectura Sistemului}

\section{Viziunea Arhitecturală}

Proiectarea oricărui sistem software complex pornește de la un set de principii
fundamentale care ghidează toate deciziile tehnice ulterioare. În cazul sistemului
descris în prezenta lucrare, trei principii arhitecturale au constituit pilonii
conceptuali ai întregului demers: localitatea datelor (\textit{local-first}),
confidențialitatea prin proiectare (\textit{privacy by design}) și inteligența la
marginea rețelei (\textit{edge intelligence}). Înțelegerea acestor principii este
esențială pentru a aprecia rațiunea din spatele fiecărei decizii tehnologice prezentate
în capitolele ulterioare.

\subsection{Principiul Local-First}

Conceptul de \textit{local-first software}, formalizat în literatura academică de
Kleppmann et al. \cite{kleppmann2019}, descrie o paradigmă de proiectare
în care datele sunt stocate și procesate în primul rând pe dispozitivele utilizatorului,
nu pe servere externe. Aplicat în contextul IoT și al sistemelor smart home, principiul
local-first presupune că toate funcționalitățile esențiale ale sistemului ---
colectarea datelor, stocarea, procesarea și controlul actuatoarelor --- trebuie să
funcționeze în mod complet și corect în absența oricărei conexiuni la internet.

Din punct de vedere tehnic, implementarea principiului local-first în sistemul de
față s-a materializat prin mai multe decizii de proiectare concrete. Broker-ul MQTT
Mosquitto rulează pe Raspberry Pi 3B+, eliminând orice intermediar cloud în fluxul de
comunicație dintre nodurile senzoriale și backend. Baza de date TimescaleDB este
găzduită local, pe același dispozitiv gateway, asigurând că întregul istoric al datelor
rămâne sub controlul exclusiv al utilizatorului. Backend-ul NestJS și algoritmul de
machine learning rulează de asemenea pe Raspberry Pi, iar aplicația mobilă comunică
direct cu adresa IP locală a gateway-ului, fără a tranzita niciun server extern.

Această abordare prezintă avantaje măsurabile față de modelul cloud-centric:
latența end-to-end pentru o alertă generată de un senzor până la notificarea pe
telefon se reduce la câteva zeci de milisecunde pe rețeaua locală, față de sute de
milisecunde sau chiar secunde specifice rutelor prin cloud. Disponibilitatea
sistemului devine independentă de furnizori externi sau de calitatea conexiunii la
internet, aspect cu implicații practice semnificative în scenarii reale.

\subsection{Confidențialitatea prin Proiectare}

\textit{Privacy by design} este un principiu formulat de Ann Cavoukian în anii '90
\cite{cavoukian2009}, ulterior adoptat ca cerință explicită de GDPR (Art. 25), care
stipulează că protecția datelor personale trebuie integrată în arhitectura sistemelor
informatice încă din faza de proiectare, nu adăugată post-factum ca strat separat.
Aplicat sistemului smart home, acest principiu are implicații directe și concrete.

Datele colectate de senzorii ambientali --- temperaturi, mișcări, niveluri de gaze,
imagini video --- constituie date personale sensibile, deoarece permit inferirea
comportamentului, rutinelor și prezenței persoanelor în locuință. În arhitectura
propusă, aceste date nu traversează niciodată rețele publice și nu sunt accesibile
unor terțe părți. Toate comunicațiile dintre nodurile ESP32 și gateway sunt protejate
prin MQTT over TLS (portul 8883), prevenind interceptarea în cadrul rețelei locale.
Autentificarea la nivelul API-ului NestJS se realizează prin token-uri JWT cu durată
limitată, împiedicând accesul neautorizat. Aplicația mobilă nu include niciun SDK
de analytics, advertising sau telemetrie terță parte.

\subsection{Inteligența la Marginea Rețelei față de Inteligența în Cloud}

Dihotomia \textit{edge intelligence} versus \textit{cloud intelligence} reprezintă
una dintre dezbaterile centrale ale IoT-ului modern. Abordarea tradițională constă
în transmiterea datelor brute de la senzori către platforme cloud (AWS IoT, Google
Cloud IoT, Azure IoT Hub) unde sunt procesate de modele ML cu resurse computaționale
nelimitate. Abordarea \textit{edge} presupune rularea algoritmilor de ML pe dispozitive
cu resurse limitate, amplasate aproape de sursa datelor.

Compromisul fundamental constă în faptul că modelele cloud beneficiază de putere
de calcul practic nelimitată, permițând utilizarea unor arhitecturi complexe (rețele
neuronale profunde, modele de mari dimensiuni), în timp ce modelele edge sunt
constrânse de resursele hardware disponibile, dar beneficiază de latență redusă,
funcționare offline și confidențialitate. Pentru cazul de utilizare specific al
sistemului de față --- detecția anomaliilor comportamentale pe date multidimensionale
de la 4--5 senzori --- algoritmul Isolation Forest reprezintă un compromis optim:
are complexitate computațională moderată, adecvată pentru Raspberry Pi 3B+, și
oferă performanțe excelente pe probleme de detecție a anomaliilor în date cu
distribuție necunoscută, fără a necesita date etichetate pentru antrenare.

\section{Arhitectura pe Straturi}

Sistemul este structurat pe cinci straturi funcționale distincte, fiecare cu
responsabilități bine delimitate și interfețe clar definite față de straturile adiacente.
Această organizare verticală respectă principiul separării preocupărilor
(\textit{separation of concerns}) și facilitează mentenanța, testarea și evoluția
independentă a fiecărei componente.

Figura~\ref{fig:arhitectura} ilustrează schema arhitecturii pe straturi a sistemului.

\begin{figure}[H]
\begin{center}
\begin{BVerbatim}[fontsize=\small]
+---------------------------------------------+
|  LAYER 5 - Aplicatie Mobila                 |
|  React Native + Expo (iOS / Android)        |
+------------------+--------------------------+
                   |  HTTP REST / WebSocket
+------------------v--------------------------+
|  LAYER 4 - Backend API                      |
|  NestJS + TypeORM + Socket.io + JWT         |
+------------------+--------------------------+
                   |  intern (acelasi device)
+------------------v--------------------------+
|  LAYER 3 - Gateway Raspberry Pi 3B+         |
|  Mosquitto MQTT  |  TimescaleDB  |  ML      |
+------------------+--------------------------+
                   |  MQTT over TLS (port 8883)
+------------------v--------------------------+
|  LAYER 2 - Noduri ESP32                     |
|  ESP32 WROOM-32 (senzori + servo)           |
|  ESP32-CAM (supraveghere video)             |
+------------------+--------------------------+
                   |  GPIO / I2C / PWM
+------------------v--------------------------+
|  LAYER 1 - Senzori & Actuatori              |
|  DHT22 | MQ-2 | PIR | BH1750 | Servo SG90  |
+---------------------------------------------+
\end{BVerbatim}
\end{center}
\caption{Arhitectura pe cinci straturi a sistemului Smart Home Local-First}
\label{fig:arhitectura}
\end{figure}

\subsection{Stratul 1 --- Senzori și Actuatori}

Primul strat al arhitecturii cuprinde totalitatea dispozitivelor fizice de intrare și
ieșire conectate la nodurile microcontroler. Din perspectiva senzorilor, sistemul
integrează patru tipuri distincte de traductoare, acoperind principalele dimensiuni
de monitorizare ale unui mediu rezidențial.

Senzorul DHT22 asigură măsurarea temperaturii în intervalul $-40\,^{\circ}C \ldots +80\,^{\circ}C$
cu o precizie de $\pm 0{,}5\,^{\circ}C$ și a umidității relative în intervalul
$0\,\%\ldots 100\,\%$ cu o precizie de $\pm 2$--$5\,\%$. Comunicația se realizează
prin protocolul 1-Wire proprietar pe un singur pin digital. Senzorul MQ-2 este un
detector de gaze combustibile și fum, bazat pe un element de oxid metalic semiconductor
(SnO\textsubscript{2}) al cărui rezistivitate se modifică în prezența gazelor
reducătoare (metan, propan, hidrogen, monoxid de carbon, fum). Semnalul de ieșire este
analogic, necesitând conversie analog-digitală prin ADC-ul intern al ESP32. Senzorul
PIR (Passive Infrared) HC-SR501 detectează mișcarea prin variațiile de radiație
infraroșie generate de deplasarea corpurilor calde în câmpul său de vedere (unghi de
$120^{\circ}$, rază de 7 m), furnizând un semnal digital binar. Senzorul BH1750
măsoară luminanța ambientală în intervalul 1--65535 lux prin comunicație I\textsuperscript{2}C,
oferind date cantitative pentru automatizarea iluminatului.

Din perspectiva actuatoarelor, sistemul include un releu de putere pentru controlul
echipamentelor electrice de curent alternativ (până la 10A/250VAC), un servo motor
pentru aplicații de control mecanic (închidere/deschidere perdele), și un buzzer
piezoelectric pentru alertare acustică locală.

\subsection{Stratul 2 --- Noduri ESP32}

Al doilea strat este constituit din microcontrolerele ESP32, care joacă rolul de
interfețe inteligente între lumea fizică a senzorilor și rețeaua locală. Modulul
ESP32 WROOM-32, produs de Espressif Systems, integrează un procesor dual-core
Xtensa LX6 la 240 MHz, 4 MB de memorie flash, 520 KB SRAM, conectivitate Wi-Fi
802.11 b/g/n și Bluetooth 4.2/BLE. Aceste specificații tehnice îl poziționează
semnificativ deasupra microcontrolerelor clasice de tip Arduino, oferind resurse
computaționale suficiente pentru rularea unui stivă de protocoale complete (TCP/IP,
TLS, MQTT) și pentru preprocesarea datelor la nivel local.

Firmware-ul rulat pe nodurile ESP32 WROOM-32 este responsabil pentru: inițializarea
și calibrarea senzorilor la pornire, citirea periodică a valorilor (la interval
configurabil, implicit 30 de secunde), validarea și filtrarea datelor aberante,
serializarea valorilor în format JSON și publicarea acestora pe topicurile MQTT
corespunzătoare, precum și pentru abonarea la topicurile de comandă și executarea
acțiunilor solicitate (activare releu, poziționare servo).

Nodul ESP32-CAM reprezintă o variantă specializată a familiei ESP32, echipată
suplimentar cu un modul cameră OV2640 (2 MP, până la 1600$\times$1200 pixeli) și un
slot microSD pentru stocarea locală a imaginilor. Acesta operează independent față de
nodul senzorial principal, gestionând captarea de imagini la cerere sau la detectarea
mișcării, comprimarea JPEG și transmiterea unui flux video continuu MJPEG către
aplicația mobilă.

\subsection{Stratul 3 --- Gateway Raspberry Pi 3B+}

Raspberry Pi 3B+ constituie nucleul arhitectural al întregului sistem, concentrând
patru subsisteme funcționale critice pe un singur dispozitiv cu consum energetic
redus (5V/2,5A, aproximativ 5--7W în funcționare normală).

\textbf{Broker-ul MQTT Mosquitto} gestionează întreaga comunicație publish-subscribe
între nodurile ESP32 și backend-ul NestJS. Mosquitto este un broker open-source de
înaltă performanță, implementând standardul MQTT 3.1.1 și 5.0, capabil să gestioneze
mii de conexiuni simultane pe hardware modest. Configurarea cu TLS mutual asigură
criptarea end-to-end a tuturor mesajelor IoT. Topicurile MQTT sunt organizate ierarhic
după schema \texttt{home/\{room\}/\{sensor\_type\}} pentru date și
\texttt{home/\{room\}/cmd/\{actuator\}} pentru comenzi.

\textbf{TimescaleDB} este extensia PostgreSQL optimizată pentru date de tip serie
temporală, selectată pentru stocarea istoricului complet al citirilor senzorilor.
TimescaleDB introduce conceptul de \textit{hypertable} --- o tabelă particionată
automat pe dimensiunea temporală --- care menține performanțe constante de inserare și
interogare chiar și pe seturi de date cu milioane de înregistrări. Funcțiile de
\textit{continuous aggregation} permit calculul eficient al statisticilor agregate
(medii orare, zilnice, săptămânale) fără a reciti datele brute.

\textbf{ML Engine-ul} reprezintă componenta de inteligență artificială a sistemului,
implementată ca un serviciu Python care rulează Isolation Forest din biblioteca
\texttt{scikit-learn}. Serviciul consumă datele stocate în TimescaleDB pentru antrenare,
evaluează periodic normalitatea citirilor curente și publică rezultatele evaluării
(scor de anomalie, clasificare normal/anormal) pe topicul MQTT dedicat alertelor.

\subsection{Stratul 4 --- Backend API NestJS}

Backend-ul NestJS reprezintă stratul de orchestrare și abstractizare al sistemului,
expunând interfețe standardizate consumabile de aplicația mobilă. Ales pentru
arhitectura sa modulară și orientată pe decoratori, NestJS organizează funcționalitățile
în module independente: modulul de autentificare (JWT, bcrypt pentru hashing parole),
modulul de senzori (CRUD pentru configurații și citiri), modulul de actuatori (comenzi
și stare), modulul de alerte (stocare, marcare ca citit, filtrare) și gateway-ul
WebSocket (Socket.io, broadcast selectiv per client autentificat).

Comunicarea cu broker-ul MQTT se realizează intern, în cadrul aceluiași dispozitiv
Raspberry Pi, prin biblioteca \texttt{mqtt.js}, eliminând latența de rețea pentru
acest flux critic. Comunicarea cu TimescaleDB se realizează prin TypeORM, cu entități
mapate pe structura hypertable-urilor.

\subsection{Stratul 5 --- Aplicația Mobilă React Native/Expo}

Stratul de prezentare al sistemului este reprezentat de o aplicație mobilă
cross-platform, dezvoltată cu React Native și Expo, compatibilă cu platformele iOS
și Android. Aplicația adoptă o arhitectură bazată pe Zustand pentru managementul stării
globale și React Navigation pentru navigarea între ecrane. Comunicarea cu backend-ul
utilizează Axios pentru cererile HTTP REST și Socket.io-client pentru conexiunea
WebSocket persistentă, prin care sunt primite actualizările în timp real ale datelor
senzorilor și notificările de alertă. Interfața vizuală adoptă o temă întunecată caldă
(\textit{warm dark}), construită pe paleta de culori stone și amber, optimizată pentru
utilizare în condiții de luminozitate redusă.

\section{Fluxul de Date (End-to-End)}

Înțelegerea completă a arhitecturii necesită urmărirea datelor prin toate straturile
sistemului, atât pe calea de monitorizare (de la senzori la ecranul utilizatorului),
cât și pe calea de comandă (de la interfața mobilă la actuatori) și pe calea de alertare
(de la detectarea anomaliei la notificarea utilizatorului).

\subsection{Fluxul de Monitorizare (Senzori $\rightarrow$ Mobil)}

Ciclul de monitorizare începe cu citirea periodică a valorilor senzorilor de către
firmware-ul ESP32, la intervalele configurate. Valorile sunt validate local (eliminarea
citirilor ieșite din domeniu sau a erorilor de comunicație), serializate în format JSON
și publicate pe topicul MQTT corespunzător (de exemplu, \texttt{home/living/temperature}).
Broker-ul Mosquitto recepționează mesajul și îl livrează backend-ului NestJS, care
persistează înregistrarea în TimescaleDB, actualizează cache-ul intern cu ultima valoare
și emite un eveniment WebSocket pe canalul \texttt{sensor-update} către toți clienții
mobili autentificați. Întregul flux se desfășoară, în condiții normale de rețea locală,
în mai puțin de 100 de milisecunde de la momentul citirii senzorului.

\subsection{Fluxul de Comandă (Mobil $\rightarrow$ Actuator)}

Calea de comandă funcționează în sens invers față de fluxul de monitorizare.
Utilizatorul interacționează cu un control din aplicația mobilă (de exemplu, un
comutator pentru releul de putere). Aplicația emite o cerere HTTP POST către
endpoint-ul REST corespunzător al backend-ului NestJS, cu payload-ul comenzii și
token-ul JWT de autentificare în header-ul Authorization. Backend-ul validează
token-ul, persistă comanda și publică un mesaj MQTT pe topicul de comandă al
actuatorului (\texttt{home/living/cmd/relay}). Nodul ESP32, abonat la acest topic,
recepționează comanda și execută acțiunea fizică corespunzătoare. Starea actuatorului
este publicată înapoi pe un topic de confirmare, închizând bucla de feedback.

\subsection{Fluxul de Alertare (Anomalie ML $\rightarrow$ Notificare)}

Calea de alertare este declanșată de ML Engine-ul care rulează periodic pe Raspberry Pi.
La fiecare ciclu de evaluare, serviciul Python extrage ultimele citiri ale senzorilor din
TimescaleDB și le supune scoringului Isolation Forest. Dacă scorul de anomalie depășește
pragul calibrat, ML Engine-ul publică un mesaj MQTT pe topicul \texttt{home/alerts/anomaly},
cu detalii despre tipul anomaliei, valorile implicate și scorul de încredere. Backend-ul
NestJS persistă alerta și emite imediat un eveniment WebSocket de tipul \texttt{alert-new}
către toți clienții mobili conectați. Aplicația mobilă afișează o notificare locală și
actualizează ecranul de alerte cu noua înregistrare.

\section{Alegerea Tehnologiilor}

Fiecare decizie tehnologică din cadrul proiectului a fost fundamentată pe o analiză
comparativă a alternativelor disponibile, luând în considerare criteriile specifice
ale unui sistem local-first pentru IoT rezidențial.

\subsection{NestJS față de alternativele backend}

Alternativele evaluate pentru stratul de backend au inclus Flask și Django (Python),
Express.js (Node.js minimal) și Spring Boot (Java). NestJS a fost ales din mai multe
motive convergente. Arhitectura sa modulară și sistemul de injecție de dependențe reduc
semnificativ complexitatea accidentală a proiectului comparativ cu Express.js. Utilizarea
TypeScript ca limbaj nativ oferă siguranță de tip (\textit{type safety}) end-to-end,
reducând clasa de erori runtime și facilitând refactorizarea. Suportul nativ pentru
WebSocket prin decoratori și gateway-uri Socket.io elimină necesitatea configurării
manuale a unui server WebSocket separat. Față de soluțiile Python (Django/Flask), NestJS
oferă performanțe superioare la gestionarea conexiunilor WebSocket concurente, aspect
critic pentru un sistem cu actualizări frecvente în timp real.

\subsection{React Native față de Flutter}

Principalele alternative evaluate pentru aplicația mobilă au fost Flutter (Google, Dart)
și React Native (Meta, JavaScript/TypeScript). Alegerea React Native/Expo s-a bazat pe:
ecosistemul JavaScript extrem de matur și bogat în librării de componentă, alinierea cu
backend-ul TypeScript pentru partajarea de tipuri și logică de validare, instrumentul
Expo Go care permite iterarea rapidă fără build nativ la fiecare modificare, și
compatibilitatea bibliotecilor de grafice cu cerințele de vizualizare ale proiectului.
Flutter ar fi oferit performanțe de randare superioare prin compilarea nativă, însă costul
costul învățării limbajului Dart și ecosistemul mai puțin matur pentru librăriile IoT-specifice
au constituit dezavantaje decisive pentru contextul acestui proiect.

\subsection{MQTT față de HTTP polling}

Alegerea MQTT ca protocol de comunicație IoT a fost facilă față de alternativa HTTP polling.
Modelul publish-subscribe al MQTT este natural adaptat pentru scenariile IoT unde mulți
producători de date (noduri senzoriale) comunică asincron cu mulți consumatori (backend,
aplicații de monitorizare). Latența MQTT este semnificativ mai mică față de HTTP polling,
deoarece mesajele sunt livrate imediat la publicare. Consumul de bandă și de energie este
minim, dat fiind că headerele MQTT sunt de ordinul a zeci de bytes. Nivelurile de QoS
(Quality of Service) oferite de protocol --- QoS 0 (cel mult o dată), QoS 1 (cel puțin o dată),
QoS 2 (exact o dată) --- permit calibrarea compromisului între fiabilitate și overhead
în funcție de tipul datelor.

\subsection{TimescaleDB față de PostgreSQL simplu}

TimescaleDB extinde PostgreSQL cu optimizări specifice datelor de tip serie temporală
fără a sacrifica compatibilitatea cu ecosistemul SQL standard. Față de un PostgreSQL simplu,
TimescaleDB oferă: partiționarea automată a datelor pe dimensiunea temporală prin
hypertable-uri, cu menținerea performanțelor de inserare și interogare la scară; funcții
de agregate continue (\textit{continuous aggregations}) care calculează și mențin automat
statistici pre-agregate pentru interogări rapide pe intervale de timp; politici de compresie
automată a datelor istorice cu rapoarte de compresie de 10--20$\times$ față de tabelele
necomprimate. Față de soluții dedicate de time-series cum ar fi InfluxDB sau Prometheus,
TimescaleDB oferă avantajul compatibilității complete SQL și al utilizării ecosistemului
PostgreSQL (TypeORM, pgAdmin), reducând complexitatea infrastructurii.

\subsection{Isolation Forest față de reguli cu praguri fixe}

Sistemele clasice de alertare în smart home utilizează reguli cu praguri fixe: dacă
temperatura depășește $35\,^{\circ}C$, se generează alertă; dacă concentrația MQ-2
depășește un prag absolut, se generează alertă. Această abordare este simplă și
transparentă, dar prezintă limitări fundamentale: pragurile optime variază între locuințe,
anotimpuri și obiceiuri ale utilizatorului; nu captează corelațiile dintre multiple
dimensiuni (de exemplu, o temperatură de $30\,^{\circ}C$ este normală în august dar
anormală în ianuarie); și tind să genereze rate ridicate de fals-pozitive sau fals-negative
dacă sunt calibrate incorect.

Algoritmul Isolation Forest, propus de Liu et al. \cite{liu2008},
abordează diferit problema detecției anomaliilor: în loc să definească normalul și să
identifice devierile de la el, izolează explicit punctele anomale prin partiționarea
aleatorie recursivă a spațiului de date. Punctele anomale, fiind rare și diferite de masa
datelor normale, sunt izolate cu un număr mai mic de partiționări, obținând un scor de
anomalie ridicat. Avantajele pentru cazul de față sunt: algoritmul nu necesită date
etichetate pentru antrenare (învățare nesupervizată); are complexitate computațională
$O(n \log n)$ și consum de memorie $O(n)$, adecvat pentru Raspberry Pi; se adaptează
automat la tiparele specifice ale utilizatorului pe măsura acumulării datelor de antrenare.

\section{Considerații de Securitate}

Securitatea unui sistem smart home este o preocupare de prim ordin, dat fiind că
dispozitivele controlate (releu de putere, actuatori mecanici) pot influența direct
siguranța fizică a locuinței. Deși sistemul propus operează exclusiv pe rețeaua locală
--- eliminând o categorie întreagă de vectori de atac specifici serviciilor cloud ---
există totuși amenințări relevante care trebuie adresate prin mecanisme tehnice adecvate.

\textbf{Autentificarea și autorizarea} la nivelul API-ului NestJS se realizează prin
JSON Web Tokens (JWT) cu o durată de expirare de șapte zile. Token-urile sunt semnate
cu un secret configurat ca variabilă de mediu, prevenind falsificarea. Parola
utilizatorului este stocată exclusiv sub formă de hash bcrypt (factor de cost 12),
astfel încât chiar și în cazul unui acces neautorizat la baza de date, parola în clar
nu poate fi recuperată în timp util. Toate endpoint-urile API, cu excepția rutei de
login, sunt protejate prin middleware-ul de autentificare JWT.

\textbf{Criptarea comunicațiilor MQTT} este asigurată prin TLS (Transport Layer Security)
pe portul 8883. Certificatele TLS sunt generate local și stocate pe Raspberry Pi, iar
nodurile ESP32 sunt configurate cu certificatul CA corespunzător. Această măsură
previne interceptarea sau modificarea mesajelor MQTT de către un potențial atacator
cu acces la rețeaua Wi-Fi locală.

\textbf{Izolarea rețelei} constituie primul și cel mai important strat de securitate al
sistemului. Prin design, backend-ul NestJS expune portul API exclusiv pe interfața de
rețea locală (LAN), fără port-forwarding sau expunere directă pe internet. Utilizatorul
care dorește acces de la distanță trebuie să configureze explicit o soluție de VPN
(recomandat WireGuard), menținând controlul complet al suprafeței de atac.

Sistemul nu implementează autentificare la nivelul MQTT (autentificare per-client cu
username/parolă), aspect identificat ca limitare actuală și direcție de îmbunătățire
pentru versiunile viitoare. De asemenea, posibilitatea accesului de la distanță prin VPN,
deși menționată, nu face obiectul implementării prezente, focalizată pe scenariul de
utilizare locală.


%% ============================================================
%%  CAPITOLUL 3 — IMPLEMENTAREA HARDWARE
%% ============================================================
\chapter{Implementarea Hardware}

\section{Prezentare Generală}

Arhitectura fizică a sistemului smart home descris în această lucrare este organizată
pe două niveluri ierarhice distincte: nivelul de percepție și acțiune, reprezentat de
nodurile embedded bazate pe microcontrolere ESP32, și nivelul de procesare și
orchestrare, asigurat de un calculator monoplacă Raspberry Pi 3B+.

La nivelul de percepție funcționează doi noduri independente. Primul nod --- construit
pe microcontroler ESP32 WROOM-32 --- are rolul de a colecta date de la senzori de mediu
și de a controla actuatorii din camera de zi. Al doilea nod --- realizat pe modulul
ESP32-CAM AI-Thinker --- asigură supravegherea vizuală a intrării și holului, furnizând
un flux video continuu accesibil din aplicația mobilă.

La nivelul de procesare, Raspberry Pi 3B+ îndeplinește rolul de gateway local și
server de aplicații. Acesta găzduiește broker-ul MQTT, serviciul NestJS cu API REST
și WebSocket, baza de date PostgreSQL și motorul de detecție a anomaliilor bazat pe
Isolation Forest. Toate comunicațiile între noduri și server se desfășoară exclusiv în
rețeaua locală Wi-Fi, în conformitate cu principiul local-first enunțat în capitolul anterior.

Prototipul prezentat în această lucrare este realizat pe breadboard, configurație adecvată
fazei de cercetare și validare funcțională. O eventuală trecere în producție ar implica
proiectarea și fabricarea unui PCB dedicat, aspect care depășește sfera prezentei lucrări.

\section{Microcontrolere Utilizate}

\subsection{ESP32 WROOM-32 --- Nodul de Senzori}

ESP32 WROOM-32 este un modul de tip System-on-Module (SoM) produs de Espressif Systems,
proiectat pentru aplicații IoT care necesită conectivitate wireless și putere de calcul
ridicată la un cost redus. Nucleul modulului este cipul ESP32-D0WDQ6, cu arhitectură
dual-core Xtensa LX6, frecvență maximă de operare de 240 MHz și memorie flash internă
de 4 MB. Memoria SRAM disponibilă este de 520 KB, suficientă pentru stocarea programului
firmware și a structurilor de date asociate senzorilor.

Din punct de vedere al conectivității wireless, modulul integrează Wi-Fi 802.11 b/g/n
(2,4 GHz) și Bluetooth 4.2 (clasic și BLE), ambele pe o antenă PCB integrată. Pentru
această aplicație se utilizează exclusiv interfața Wi-Fi, prin care nodul publică datele
de telemetrie către broker-ul MQTT de pe Raspberry Pi.

Perifericele disponibile sunt deosebit de versatile: 30 de pini GPIO configurabili, un
convertor analog-digital (ADC) de 12 biți cu 18 canale, interfețe I\textsuperscript{2}C,
SPI, UART și capacitate de generare PWM pe aproape orice pin. Tensiunea de operare este
de 3,3 V, însă unii pini tolerează semnale de 5 V ca intrare.

Alegerea acestui microcontroler a fost motivată de mai mulți factori. Documentația
extensivă și comunitatea vastă de utilizatori reduc semnificativ efortul de depanare.
ADC-ul de 12 biți este necesar pentru citirea valorilor analogice de la senzorul MQ-2
de gaz. Interfața I\textsuperscript{2}C este utilizată pentru comunicarea cu senzorul
BH1750 de lumină. Capacitatea PWM pe GPIO18 permite controlul precis al servomotorului
SG90. Nu în ultimul rând, prețul accesibil al modulului (aproximativ 5--7 EUR) contribuie
la menținerea costului total al sistemului la un nivel rezonabil.

\subsection{ESP32-CAM AI-Thinker --- Nodul Cameră}

Modulul ESP32-CAM AI-Thinker este o soluție integrată care combină un microcontroler
ESP32-S cu un modul de cameră OV2640 de 2 megapixeli, destinată aplicațiilor de
supraveghere video cu cost redus. Pe lângă procesorul ESP32-S, modulul dispune de o
memorie PSRAM (Pseudo-Static RAM) de 4 MB, necesară pentru stocarea temporară a
cadrelor video înainte de transmisie, și un slot microSD pentru înregistrare locală
opțională.

Senzorul de imagine OV2640 suportă rezoluții de până la UXGA (1600$\times$1200 pixeli)
și poate transmite fluxuri video comprimate în format MJPEG. Pentru aplicația prezentată,
rezoluția utilizată în mod curent este VGA (640$\times$480), care oferă un echilibru
acceptabil între calitatea imaginii și lățimea de bandă consumată în rețeaua locală.
Modulul dispune și de un LED flash pe GPIO4, utilizabil pentru iluminarea scenei în
condiții de lumină redusă.

O particularitate importantă a acestui modul este absența unui convertor USB-UART
integrat pe placă. În consecință, programarea firmware-ului necesită utilizarea unui
adaptor extern FTDI (sau similar), procedură detaliată în secțiunea~\ref{sec:espcam-prog}.

Motivele alegerii ESP32-CAM sunt în principal economice și funcționale: modulul oferă
integrare completă cameră + Wi-Fi + microcontroler la un preț de aproximativ 8 EUR,
biblioteca de streaming MJPEG este disponibilă nativ în framework-ul Arduino pentru
ESP32, iar dimensiunea compactă permite montarea discretă la intrarea locuinței.

\section{Senzori Utilizați}

\subsection{DHT22 --- Temperatură și Umiditate}

Senzorul DHT22 (cunoscut și sub denumirea AM2302) este un senzor digital de temperatură
și umiditate relativă. Domeniul de măsurare a temperaturii este cuprins între
$-40\,^{\circ}C$ și $+80\,^{\circ}C$, cu o precizie de $\pm 0{,}5\,^{\circ}C$, iar
umiditatea relativă poate fi citită în intervalul 0--100\% cu o precizie de
$\pm 2$--$5\,\%$.

Comunicarea cu microcontrolerul se realizează printr-un protocol proprietar one-wire
pe un singur pin de date. Linia DATA necesită o rezistență pull-up de 10 k$\Omega$ la
tensiunea de alimentare (3,3 V) pentru funcționare corectă. Intervalul minim între două
citiri succesive este de 2 secunde, limitat de procesul intern de conversie a senzorului.
Senzorul este conectat la GPIO4 al ESP32 WROOM-32 și alimentat de la pinul 3,3 V al
microcontrolerului.

\subsection{MQ-2 --- Gaz și Fum}

Senzorul MQ-2 este un detector electrochimic capabil să detecteze o gamă largă de gaze
inflamabile și vapori: GLP (gaz lichefiat petrolier), butan, propan, metan, hidrogen,
fum și monoxid de carbon (CO). Modulul furnizează atât un semnal analogic (0--3,3 V,
corespunzând unui interval ADC de 0--1023 unități) cât și un semnal digital de prag,
reglabil prin potențiometrul integrat pe modul.

Senzorul necesită alimentare la 5 V (curentul de consum al elementului sensibil
depășind capacitatea regulatorului de 3,3 V al ESP32) și un timp de preîncălzire de
aproximativ 20 de secunde după pornire, înainte ca valorile citite să fie stabile și
reprezentative. Pragul de alertă utilizat în firmware este de 350 de unități ADC,
valoare determinată empiric în condiții normale de ambient. Semnalul analogic este citit
de pe GPIO34 (canalul ADC1\_CH6), iar semnalul digital de prag de pe GPIO35 --- ambii
pini configurați exclusiv ca intrări.

\subsection{HC-SR501 --- Mișcare PIR}

Senzorul de mișcare HC-SR501 funcționează pe principiul detecției infraroșului pasiv
(PIR --- \textit{Passive InfraRed}). Elementul sensibil detectează variațiile de radiație
infraroșie termică emise de corpuri umane în mișcare, fără a emite nicio radiație proprie.
Raza de detecție este de până la 7 metri, cu un unghi de acoperire de $120^{\circ}$.

Modulul se alimentează la 5 V și furnizează un semnal digital de ieșire de 3,3 V
(compatibil direct cu intrările ESP32), activ HIGH pe durata detecției de mișcare.
Sensibilitatea detectivă și durata menținerii semnalului HIGH după detecție sunt reglabile
prin intermediul a doi potențiometre aflați pe spatele modulului. Senzorul este conectat
la GPIO27 al ESP32.

\subsection{BH1750 --- Lumină Ambientală}

BH1750 este un senzor digital de iluminanță (luminos ambientală) cu rezoluție de
16 biți, capabil să măsoare niveluri de lumină în intervalul 1--65535 lux. Comunicarea
cu microcontrolerul se realizează prin interfața I\textsuperscript{2}C, la adresa 0x23
atunci când pinul ADDR este conectat la GND (adresa alternativă 0x5C se obține prin
conectarea ADDR la VCC).

Senzorul se alimentează la 3,3 V, ceea ce îl face compatibil direct cu pinii ESP32.
Pinii I\textsuperscript{2}C utilizați sunt GPIO21 (SDA) și GPIO22 (SCL), care reprezintă
pinii I\textsuperscript{2}C impliciți ai ESP32. Datele de iluminanță sunt utilizate atât
pentru logica de automatizare a iluminatului (control releu și servo perdele) cât și ca
\textit{feature} în algoritmul de detecție a anomaliilor comportamentale.

\section{Actuatori Utilizați}

\subsection{Modulul Releu 4 Canale}

Modulul de releu cu 4 canale este un circuit de comandă care permite controlul sarcinilor
electrice de tensiune și curent ridicate (230 VAC, 50 Hz în rețeaua electrică de uz
casnic) prin intermediul semnalelor logice de joasă tensiune furnizate de microcontroler.
Fiecare canal este izolat optic prin optocuplori, protejând astfel microcontrolerul de
eventualele perturbații electrice generate de sarcinile comutate.

Releele sunt de tip activ LOW: starea activată (contact normal-deschis închis) corespunde
unui semnal LOW (0 V) pe pinul de intrare. Acest comportament trebuie luat în considerare
la inițializarea firmware-ului pentru a evita activarea accidentală a releelor la pornire.
Bobina fiecărui releu necesită 5 V și aproximativ 70--80 mA, motiv pentru care modulul
este alimentat dintr-o sursă externă și nu direct din pinii GPIO ale ESP32. Capacitatea
maximă de comutare este de 10 A la 250 VAC sau 10 A la 30 VDC per canal. Pinii de
control sunt GPIO26, GPIO25, GPIO33 și GPIO32.

\subsection{Buzzer Activ}

Buzzer-ul activ este un transductor electroacustic cu oscilator intern integrat, care
produce un semnal sonor continuu atunci când i se aplică tensiune de alimentare (5 V),
fără a necesita un semnal PWM extern. Simplitatea de control --- activare/dezactivare
prin nivel logic pe GPIO5 --- îl face adecvat pentru alerte locale de urgență, în
special în scenariile de detecție a gazului sau tentativă de efracție.

\subsection{Servomotorul SG90 --- Control Perdele}

SG90 este un micro-servomotor cu angrenaje din plastic, utilizat pe scară largă în
aplicații robotice și de automatizare la scară mică. Oferă un cuplu de 1,8 kg$\cdot$cm
la 5 V și poate efectua rotații în intervalul $0^{\circ}$--$180^{\circ}$. Controlul
poziției se realizează prin semnal PWM la frecvența de 50 Hz: lățimea impulsului variază
de la 0,5 ms (corespunzând poziției de $0^{\circ}$) la 2,5 ms (corespunzând poziției
de $180^{\circ}$).

În cadrul sistemului, servomotorul este utilizat pentru acționarea unui mecanism simplu
de deschidere/închidere a perdelelor, cu trei poziții predefinite: $0^{\circ}$ (perdele
complet închise), $90^{\circ}$ (perdele la jumătate) și $180^{\circ}$ (perdele complet
deschise). Aceasta permite automatizarea iluminatului natural în funcție de nivelul de
lumină ambientală măsurat de senzorul BH1750.

Alimentarea servomotorului trebuie realizată obligatoriu dintr-o sursă externă de 5 V,
și nu din pinii de alimentare ai ESP32. Un servomotor SG90 poate consuma curenți de
vârf de 500--600 mA în condiții de sarcină, valoare ce depășește cu mult capacitatea
de furnizare de curent a regulatorului intern al ESP32 (maxim 40 mA per pin GPIO,
respectiv câteva sute de mA pe pinul 3V3). Semnalul PWM este generat pe GPIO18, pin
cu capacitate PWM nativă.

\section{Schema de Conexiuni}
\label{sec:schema}

Tabelul~\ref{tab:conexiuni} prezintă complet conexiunile dintre componentele hardware
și pinii corespunzători ai ESP32 WROOM-32, împreună cu observațiile relevante pentru
asamblarea corectă a prototipului pe breadboard.

\begin{table}[H]
\caption{Conexiunile componentelor pe nodul ESP32 WROOM-32}
\label{tab:conexiuni}
\centering
\begin{tabular}{llll}
\toprule
\textbf{Componentă} & \textbf{Pin componentă} & \textbf{Pin ESP32} & \textbf{Obs.} \\
\midrule
DHT22           & VCC       & 3,3V (intern)  & ---                    \\
DHT22           & GND       & GND            & ---                    \\
DHT22           & DATA      & GPIO4          & 10 k$\Omega$ pull-up la 3,3V \\
MQ-2            & VCC       & 5V (extern)    & Nu din ESP32!          \\
MQ-2            & GND       & GND            & ---                    \\
MQ-2            & A0        & GPIO34         & ADC1\_CH6              \\
MQ-2            & D0        & GPIO35         & Input-only             \\
PIR HC-SR501    & VCC       & 5V (extern)    & ---                    \\
PIR HC-SR501    & GND       & GND            & ---                    \\
PIR HC-SR501    & OUT       & GPIO27         & ---                    \\
BH1750          & VCC       & 3,3V (intern)  & ---                    \\
BH1750          & GND       & GND            & ---                    \\
BH1750          & SDA       & GPIO21         & I\textsuperscript{2}C default \\
BH1750          & SCL       & GPIO22         & I\textsuperscript{2}C default \\
BH1750          & ADDR      & GND            & Adresă 0x23            \\
Relay 4ch       & VCC       & 5V (extern)    & ---                    \\
Relay 4ch       & GND       & GND            & ---                    \\
Relay 4ch       & IN1       & GPIO26         & Active LOW             \\
Relay 4ch       & IN2       & GPIO25         & Active LOW             \\
Relay 4ch       & IN3       & GPIO33         & Active LOW             \\
Relay 4ch       & IN4       & GPIO32         & Active LOW             \\
Buzzer          & +         & GPIO5          & ---                    \\
Buzzer          & --        & GND            & ---                    \\
Servo SG90      & VCC roșu  & 5V (extern!)   & Nu din ESP32!          \\
Servo SG90      & GND maro  & GND            & ---                    \\
Servo SG90      & PWM galben & GPIO18        & 50 Hz PWM              \\
\bottomrule
\end{tabular}
\end{table}

\section{Alimentarea Sistemului}

Diversitatea componentelor și cerințele lor electrice diferite impun o strategie atentă
de alimentare. Modulul de releu, senzorul MQ-2, senzorul PIR HC-SR501 și servomotorul
SG90 necesită toți tensiunea de 5 V și pot solicita curenți cumulativi de ordinul
amperilor în condiții de operare simultană. Aceste valori depășesc capacitatea de
furnizare a regulatorului intern al ESP32 sau a unui port USB standard.

Soluția adoptată este utilizarea unei surse externe de 5 V cu curentul de ieșire de
3 A, conectată direct la șinele de alimentare ale breadboard-ului. Microcontrolerul
ESP32 WROOM-32 este alimentat fie prin conectorul USB propriu, fie prin pinul VIN
(care acceptă 5 V direct), extras din aceeași șină de 5 V a breadboard-ului. Regulatorul
LDO integrat pe modulul ESP32 generează apoi tensiunea internă de 3,3 V, disponibilă
pe pinul 3V3 și utilizată pentru alimentarea senzorilor DHT22 și BH1750, care au
consum redus (sub 5 mA fiecare).

Este esențial ca masa (GND) surselor de 5 V externe și masa ESP32 să fie conectate în
comun pe breadboard, altfel referințele de tensiune diferă și semnalele logice devin
nesigure. Principiul fundamental care trebuie respectat este că servo-ul și releele nu
trebuie alimentate niciodată direct din pinii GPIO sau de pe pinul 3V3 al ESP32 ---
o astfel de conexiune riscă distrugerea permanentă a microcontrolerului prin supracurent.

\section{Nodul Cameră ESP32-CAM --- Programare și Configurare}
\label{sec:espcam-prog}

Absența unui convertor USB-UART integrat pe modulul ESP32-CAM impune utilizarea unui
programator extern pentru încărcarea firmware-ului. Procedura standard folosește un
adaptor FTDI (FT232RL sau similar), care realizează conversia USB-Serial necesară
comunicării cu portul UART0 al ESP32.

Conexiunile necesare între adaptorul FTDI și modulul ESP32-CAM sunt următoarele:
GND (FTDI) $\rightarrow$ GND (ESP32-CAM), VCC 5V (FTDI) $\rightarrow$ 5V (ESP32-CAM),
TX (FTDI) $\rightarrow$ U0RXD / GPIO3 (ESP32-CAM), RX (FTDI) $\rightarrow$ U0TXD / GPIO1
(ESP32-CAM).

Pentru intrarea în modul de programare (\textit{flash mode}), pinul GPIO0 al ESP32-CAM
trebuie conectat la GND înainte de aplicarea tensiunii de alimentare. Această configurare
determină boot-loader-ul să aștepte transferul firmware-ului prin UART în loc să execute
codul din flash. Pașii de urmat sunt: (1) se conectează GPIO0 la GND prin intermediul
unui jumper sau fir pe breadboard; (2) se conectează adaptorul FTDI la calculator prin
USB; (3) se apasă butonul RESET de pe modulul ESP32-CAM; (4) se efectuează upload-ul
firmware-ului din Arduino IDE sau PlatformIO; (5) după finalizarea încărcării, se
îndepărtează jumper-ul GPIO0-GND; (6) se apasă din nou butonul RESET pentru pornirea
normală cu noul firmware.

Din punct de vedere funcțional, nodul cameră rulează un server MJPEG pe portul 80.
Fluxul video continuu este accesibil la adresa \texttt{http://[IP]/stream}, iar capturile
individuale (snapshot) la \texttt{http://[IP]/capture}. Spre deosebire de nodul de
senzori, nodul cameră nu publică date prin MQTT și nu participă la logica de automatizare
sau detecție a anomaliilor --- el este un nod pasiv, accesat direct din aplicația mobilă
atunci când utilizatorul dorește să vizualizeze imaginea live a intrării.

\section{Considerații Practice pentru Breadboard}

Asamblarea unui prototip funcțional pe breadboard implică o serie de constrângeri
practice care, dacă nu sunt respectate, pot genera defecțiuni intermitente dificil
de diagnosticat.

Convenția de culori pentru fire este recomandată a fi respectată în mod consistent
pe tot parcursul asamblării: roșu pentru VCC (indiferent de tensiune --- 3,3 V sau 5 V),
negru pentru GND, iar culorile neutre (galben, albastru, verde, portocaliu) pentru liniile
de semnal. Această practică reduce semnificativ riscul inversării accidentale a polarității,
mai ales în condițiile unui număr mare de conexiuni pe breadboard.

Senzorul MQ-2 necesită un timp de preîncălzire de aproximativ 20 de secunde după pornire
înainte ca rezistența elementului sensibil să se stabilizeze și valorile ADC să fie
reprezentative. Firmware-ul ignoră primele citiri după pornire pentru a evita alerte false
generate de instabilitatea de start-up a senzorului.

Senzorul DHT22 trebuie poziționat fizic la distanță de componentele care generează căldură
locală, în special bobinele releelor active. Un releu activat generează un câmp magnetic
și disipă căldură care poate influența temperatura măsurată de DHT22, introducând erori
sistematice în datele de telemetrie.

Servomotorul SG90 prezintă vibrații mecanice în timpul rotației care pot slăbi conexiunile
pe breadboard. Se recomandă fixarea firelor servomotorului cu bandă adezivă sau cleme de
fixare suplimentare pentru a preveni întreruperea intermitentă a contactelor.

Modulul ESP32-CAM se încălzește notabil în timpul transmisiei video continue, disipând
câțiva wați prin carcasa cipului. Este recomandat să se asigure un spațiu liber în jurul
modulului pentru convecție naturală a aerului, evitând acoperirea sau izolarea termică
accidentală în timpul testelor de lungă durată.


%% ============================================================
%%  CAPITOLE VIITOARE (placeholder)
%% ============================================================
\chapter{Implementarea Software}

\textit{Capitolul urmează să fie redactat. Va descrie: firmware-ul MicroPython pentru
nodurile ESP32, configurarea Mosquitto MQTT și TimescaleDB pe Raspberry Pi, arhitectura
backend-ului NestJS (module, controllere, servicii, gateway WebSocket) și implementarea
aplicației mobile React Native/Expo (ecrane, navigare, store Zustand, integrare API).}

\chapter{Detecția Anomaliilor prin Machine Learning}

\textit{Capitolul urmează să fie redactat. Va descrie: fundamente teoretice ale
algoritmului Isolation Forest (Liu et al., 2008), schema de colectare și preprocesare a
datelor de antrenare, implementarea Python cu scikit-learn pe Raspberry Pi, evaluarea
performanțelor și calibrarea pragului de anomalie pe date reale.}

\chapter{Testare și Evaluare}

\textit{Capitolul urmează să fie redactat. Va descrie: metodologia de testare
(teste funcționale, teste de performanță, teste de regresie), rezultatele testelor,
analiza consumului de resurse pe Raspberry Pi, comparație cu soluțiile comerciale
și direcții de dezvoltare ulterioară.}


%% ============================================================
%%  CONCLUZII
%% ============================================================
\chapter*{Concluzii}
\addcontentsline{toc}{chapter}{Concluzii}

\textit{Secțiunea urmează să fie completată după finalizarea tuturor capitolelor de
implementare. Va sintetiza cele mai importante contribuții ale lucrării, va evalua
gradul de îndeplinire a obiectivelor stabilite în Capitolul 1 și va propune direcții
de cercetare și dezvoltare ulterioară.}


%% ============================================================
%%  BIBLIOGRAFIE
%% ============================================================
\chapter*{Bibliografie}
\addcontentsline{toc}{chapter}{Bibliografie}

\begin{thebibliography}{99}

\bibitem{kleppmann2019}
Kleppmann, M., Wiggins, A., van Hardenberg, P., McGranaghan, M.,
\textit{Local-First Software: You Own Your Data, in Spite of the Cloud}.
ACM SIGPLAN International Symposium on New Ideas, New Paradigms, and
Reflections on Programming and Software (Onward! 2019), pp. 154--178, 2019.

\bibitem{liu2008}
Liu, F.~T., Ting, K.~M., Zhou, Z.-H.,
\textit{Isolation Forest}.
IEEE International Conference on Data Mining (ICDM), pp. 413--422, 2008.

\bibitem{cavoukian2009}
Cavoukian, A.,
\textit{Privacy by Design: The 7 Foundational Principles --- Implementation and
Mapping of Fair Information Practices}.
Information and Privacy Commissioner of Ontario, Canada, 2009.

\bibitem{espressif2023}
Espressif Systems,
\textit{ESP32 Technical Reference Manual v5.1}.
Disponibil la: \url{https://www.espressif.com/sites/default/files/documentation/esp32_technical_reference_manual_en.pdf}
(accesat: mai 2025).

\bibitem{statista2024}
*** \textit{Smart Home --- Worldwide. Market Data Report}.
Statista, 2024.
Disponibil la: \url{https://www.statista.com/outlook/dmo/smart-home/worldwide}
(accesat: mai 2025).

\bibitem{mosquitto}
*** \textit{Eclipse Mosquitto: An open source MQTT broker}.
Eclipse Foundation, 2024.
Disponibil la: \url{https://mosquitto.org/}
(accesat: mai 2025).

\bibitem{timescaledb}
*** \textit{TimescaleDB Documentation}.
Timescale, Inc., 2024.
Disponibil la: \url{https://docs.timescale.com/}
(accesat: mai 2025).

\bibitem{nestjs}
*** \textit{NestJS --- A progressive Node.js framework}.
Disponibil la: \url{https://nestjs.com/}
(accesat: mai 2025).

\bibitem{reactnative}
*** \textit{React Native --- Learn once, write anywhere}.
Meta Platforms, Inc., 2024.
Disponibil la: \url{https://reactnative.dev/}
(accesat: mai 2025).

\bibitem{expo}
*** \textit{Expo Documentation}.
Expo Technology, Inc., 2024.
Disponibil la: \url{https://docs.expo.dev/}
(accesat: mai 2025).

\end{thebibliography}


%% ============================================================
%%  DECLARAȚIE DE AUTENTICITATE
%% ============================================================
\newpage
\thispagestyle{empty}
\begin{center}
{\large\bfseries DECLARAȚIE DE AUTENTICITATE\\[4pt]
A LUCRĂRII DE FINALIZARE A STUDIILOR}
\end{center}

\vspace{8pt}

Subsemnatul \underline{\hspace{8cm}},

legitimat cu \underline{\hspace{3cm}} seria \underline{\hspace{1.5cm}}
nr. \underline{\hspace{3cm}},

CNP \underline{\hspace{8cm}},

autorul lucrării \underline{\hspace{10cm}}

\vspace{6pt}
\noindent elaborată în vederea susținerii examenului de finalizare a studiilor de
\underline{\hspace{4cm}} organizat de către Facultatea
\underline{\hspace{5cm}} din cadrul Universității Politehnica Timișoara,
sesiunea \underline{\hspace{2.5cm}} a anului universitar \underline{\hspace{2cm}},
coordonator \underline{\hspace{4cm}}, luând în considerare art. 34 din Regulamentul
privind organizarea și desfășurarea examenelor de licență/diplomă și disertație,
aprobat prin HS nr. 109/14.05.2020 și cunoscând faptul că în cazul constatării
ulterioare a unor declarații false, voi suporta sancțiunea administrativă prevăzută
de art. 146 din Legea nr. 1/2011, declar pe proprie răspundere, că:

\begin{itemize}
  \item această lucrare este rezultatul propriei activități intelectuale,
  \item lucrarea nu conține texte, date sau elemente de grafică din alte lucrări sau
        din alte surse fără ca acestea să nu fie citate,
  \item sursele bibliografice au fost folosite cu respectarea legislației române și a
        convențiilor internaționale privind drepturile de autor,
  \item această lucrare nu a mai fost prezentată în fața unei alte comisii de examen.
\end{itemize}

\vspace{1.5cm}

\noindent Timișoara,

\noindent Data \underline{\hspace{4cm}} \hfill Semnătura \underline{\hspace{5cm}}

\vspace{0.5cm}

{\footnotesize\textit{Declarația se completează „de mână" și se inserează în lucrarea
de finalizare a studiilor, la sfârșitul acesteia, ca parte integrantă.}}


%% ================================================================
\end{document}
%% ================================================================
